\section*{Chapter \ref{ch:g10:graphs-first-look} --- Graphs: Networks of Roads and People}

\begin{solution}{exo:g10:graphs-first-look:1}
The \omterm{def:g10:graphs-first-look:degree}{degrees} add up to $3 + 2 + 2 + 2 + 1 = 10$, twice the number of \omterm{def:g10:graphs-first-look:graph}{edges}: $5$ \omterm{def:g10:graphs-first-look:graph}{edges}. \omterm{def:g10:graphs-first-look:degree}{Degrees} $3$, $3$, $1$ add up to
$7$, an odd number, which cannot be twice a number of \omterm{def:g10:graphs-first-look:graph}{edges}; besides, with $3$ \omterm{def:g10:graphs-first-look:graph}{vertices} a \omterm{def:g10:graphs-first-look:degree}{degree} is at most $2$.
\end{solution}

\begin{solution}{exo:g10:graphs-first-look:2}
\code{{"A": {"B", "C"}, "B": {"A", "C", "D"}, "C": {"A", "B"}, "D": {"B"}}}: each of the four \omterm{def:g10:graphs-first-look:graph}{edges} appears twice.
\end{solution}

\begin{solution}{exo:g10:graphs-first-look:3}
A to G: A--B--C--F--G, \omterm{def:g10:graphs-first-look:distance}{distance} $4$. E to F: E--D--C--F, \omterm{def:g10:graphs-first-look:distance}{distance} $3$. From B to D there is only B--C--D: the \omterm{def:g10:graphs-first-look:graph}{graph}
has no cycle, so between two \omterm{def:g10:graphs-first-look:graph}{vertices} there is a single \omterm{def:g10:graphs-first-look:path}{path} that does not repeat a \omterm{def:g10:graphs-first-look:graph}{vertex}.
\end{solution}

\begin{solution}{exo:g10:graphs-first-look:4}
$\{1, 2, 3\}$, $\{4, 5\}$ and $\{6\}$. Each new \omterm{def:g10:graphs-first-look:graph}{edge} joins at most two components into one, so going from $3$
components to $1$ takes at least $2$ \omterm{def:g10:graphs-first-look:graph}{edges}, for example $3$--$4$ and $5$--$6$.
\end{solution}

\begin{solution}{exo:g10:graphs-first-look:5}
\code{return len(g[v])} (\code{exercises/graph_tools.py}). The people of \omterm{def:g10:graphs-first-look:degree}{degree} $1$ are $4$, $15$, $18$, $21$, $22$,
$24$, $25$, $26$, $31$ and $35$: the ten of the \omterm{def:g10:graphs-first-look:degree}{degree} table.
\end{solution}

\begin{solution}{exo:g10:graphs-first-look:6}
Each friendship is stored only at its first person, so the check of symmetry prints all $50$ friendships as \omterm{def:g10:graphs-first-look:graph}{edges}
without a reverse, instead of an empty list. The correct function also writes \code{g[b].add(a)}
(\code{exercises/graph_tools.py}).
\end{solution}

\begin{solution}{exo:g10:graphs-first-look:7}
Return \code{False} for an empty list or a first \omterm{def:g10:graphs-first-look:graph}{vertex} outside the \omterm{def:g10:graphs-first-look:graph}{graph}; otherwise check that every \omterm{def:g10:graphs-first-look:graph}{vertex} but the
last has the next one among its \omterm{def:g10:graphs-first-look:degree}{neighbours}: \code{all(seq[i + 1] in g[seq[i]] for i in range(len(seq) - 1))}
(\code{exercises/graph_tools.py}). $3$--$2$, $2$--$10$, $10$--$1$ and $1$--$28$ are friendships; $3$--$1$ is not.
\end{solution}

\begin{solution}{exo:g10:graphs-first-look:8}
Eccentricities: $1$: $2$, $2$: $2$, $3$: $3$, $4$: $2$, $5$: $3$ (the farthest pair is $3$ and $5$, through $2$ or
$4$ and then $1$). \omterm{def:g10:graphs-first-look:eccentricity}{Radius} $2$, \omterm{def:g10:graphs-first-look:eccentricity}{diameter} $3$, \omterm{def:g10:graphs-first-look:eccentricity}{centre} $\{1, 2, 4\}$.
\end{solution}

\begin{solution}{exo:g10:graphs-first-look:9}
A: in $1$, out $1$; B: in $1$, out $2$; C: in $1$, out $1$; D: in $1$, out $0$. Both sums are $4$, the number of
\omterm{def:g10:graphs-first-look:graph}{edges}: every directed \omterm{def:g10:graphs-first-look:graph}{edge} leaves exactly one \omterm{def:g10:graphs-first-look:graph}{vertex} and arrives at exactly one.
\end{solution}

\begin{solution}{exo:g10:graphs-first-look:10}
\code{return {(min(a, b), max(a, b)) for a in g for b in g[a]}}: each \omterm{def:g10:graphs-first-look:graph}{edge} is met twice, once from each end, and the
set keeps it once (\code{exercises/graph_tools.py}). The friendship \omterm{def:g10:graphs-first-look:graph}{graph} gives $50$ \omterm{def:g10:graphs-first-look:graph}{edges}, and the \omterm{def:g10:graphs-first-look:degree}{degrees} add up to
$100$.
\end{solution}

\begin{solution}{exo:g10:graphs-first-look:11}
The time is 09:30:15 (on the satellites' universal time). Latitude: $48 + 7.0381/60 \approx 48.11730$ \omterm{def:g10:graphs-first-look:degree}{degrees} north.
Longitude: $1 + 31.5240/60 = 1.52540$ \omterm{def:g10:graphs-first-look:degree}{degrees} west, so $-1.52540$ (\code{exercises/nmea.py}). The check lets the
receiving program detect a line damaged on the way: if one character changed, the exclusive or almost surely
differs from \texttt{41}, and the line is thrown away.
\end{solution}

\begin{solution}{exo:g10:graphs-first-look:12}
Count the pairs (\omterm{def:g10:graphs-first-look:graph}{vertex}, \omterm{def:g10:graphs-first-look:graph}{edge}) where the \omterm{def:g10:graphs-first-look:graph}{edge} touches the \omterm{def:g10:graphs-first-look:graph}{vertex}. Counted by \omterm{def:g10:graphs-first-look:graph}{vertex}, each \omterm{def:g10:graphs-first-look:graph}{vertex} $v$ is in
$\deg(v)$ pairs, which gives the sum of the \omterm{def:g10:graphs-first-look:degree}{degrees}. Counted by \omterm{def:g10:graphs-first-look:graph}{edge}, each \omterm{def:g10:graphs-first-look:graph}{edge} touches exactly two \omterm{def:g10:graphs-first-look:graph}{vertices}, which
gives $2|E|$. So the sum of the \omterm{def:g10:graphs-first-look:degree}{degrees} is $2|E|$, an even number; the \omterm{def:g10:graphs-first-look:graph}{vertices} of even \omterm{def:g10:graphs-first-look:degree}{degree} add an even amount,
so the odd \omterm{def:g10:graphs-first-look:degree}{degrees} must add up to an even number too, and there is an even number of them.
\end{solution}

\begin{solution}{exo:g10:graphs-first-look:13}
The \omterm{def:g10:graphs-first-look:eccentricity}{radius} is the smallest \omterm{def:g10:graphs-first-look:eccentricity}{eccentricity} and the \omterm{def:g10:graphs-first-look:eccentricity}{diameter} the largest, so $\text{radius} \le \text{diameter}$. Take
$c$ in the \omterm{def:g10:graphs-first-look:eccentricity}{centre} and any two \omterm{def:g10:graphs-first-look:graph}{vertices} $a$ and $b$. A \omterm{def:g10:graphs-first-look:path}{path} from $a$ to $c$ followed by a \omterm{def:g10:graphs-first-look:path}{path} from $c$ to $b$ is a
\omterm{def:g10:graphs-first-look:path}{path} from $a$ to $b$ (it may be shortened if it repeats a \omterm{def:g10:graphs-first-look:graph}{vertex}), so
$d(a, b) \le d(a, c) + d(c, b) \le \text{radius} + \text{radius}$. This holds for every pair, so the \omterm{def:g10:graphs-first-look:eccentricity}{diameter} is at
most $2 \times \text{radius}$. The \omterm{def:g10:graphs-first-look:graph}{graph} of eccentricities reaches it: \omterm{def:g10:graphs-first-look:eccentricity}{radius} $2$, \omterm{def:g10:graphs-first-look:eccentricity}{diameter} $4$.
\end{solution}

\begin{solution}{exo:g10:graphs-first-look:14}
\code{common_friends} returns \code{g[a] & g[b]}. A triangle $\{x, y, z\}$ is found from its \omterm{def:g10:graphs-first-look:graph}{edge} $x$--$y$ (common
friend $z$), from $x$--$z$ (common friend $y$) and from $y$--$z$ (common friend $x$): three times, so the sum is
divided by $3$ (\code{exercises/graph_tools.py}). The friendship \omterm{def:g10:graphs-first-look:graph}{graph} has $2$ triangles.
\end{solution}

\begin{solution}{exo:g10:graphs-first-look:15}
\omterm{def:g10:graphs-first-look:degree}{In-degrees}: P3 has $3$, P1, P2 and P4 have $1$, P5 has $0$. Votes: P1 gives $\tfrac12$ to P2 and $\tfrac12$ to P3;
P2, P4 and P3 give their whole vote to P3, P3 and P1; P5 gives its vote to P4. Totals: P3 $2.5$, P1 $1$, P4 $1$,
P2 $0.5$, P5 $0$. P3 is first both times. A link is a recommendation: one from a page that many others recommend
says more than one from a page nobody links to, which is why search engines weigh links by the importance of the page
they come from, and compute that importance the same way, over and over, until it settles.
\end{solution}

\begin{solution}{pb:g10:graphs-first-look:1}
\textbf{1.} $40$ \omterm{def:g10:graphs-first-look:graph}{vertices} and $50$ \omterm{def:g10:graphs-first-look:graph}{edges}; the \omterm{def:g10:graphs-first-look:degree}{degrees} add up to $100 = 2 \times 50$.

\textbf{2.} Most friends, $5$: persons $1$, $12$, $23$ and $28$. Exactly one: $4$, $15$, $18$, $21$, $22$, $24$,
$25$, $26$, $31$ and $35$.

\textbf{3.} \omterm{def:g10:graphs-first-look:degree}{Degrees} $1$, $3$ and $5$: $10 + 10 + 4 = 24$ people, an even number, as the exercise proves.

\textbf{4.} $1$, $34$, $36$, $37$ and $39$.

\textbf{5.} $29$, $10$ and $1$.

\textbf{6.} No: person $23$ is in the component of $10$ people, and no \omterm{def:g10:graphs-first-look:path}{path} joins it to the component of person $1$.

\textbf{7.} $1$--$28$, $5$--$37$, $12$--$20$ and $25$--$37$. None joins two components: $20$ and $25$ have all their
friends in the largest component.

\textbf{8.} $2$: one \omterm{def:g10:graphs-first-look:graph}{edge} joining the component of $10$ people to the largest one, and one joining person $40$.

\textbf{9.} $4$: the farthest people, $31$, $35$ and $38$, are at \omterm{def:g10:graphs-first-look:distance}{distance} $4$.

\textbf{10.} For example $3$ and $31$: $3$--$2$--$10$--$1$--$28$--$34$--$29$--$31$, seven \omterm{def:g10:graphs-first-look:graph}{edges}.

\textbf{11.} $31$, $35$ and $38$.

\textbf{12.} Person $3$ has \omterm{def:g10:graphs-first-look:eccentricity}{eccentricity} $7$: $7$ days.

\textbf{13.} \omterm{def:g10:graphs-first-look:eccentricity}{Radius} $4$, \omterm{def:g10:graphs-first-look:eccentricity}{diameter} $7$, \omterm{def:g10:graphs-first-look:eccentricity}{centre} $\{1, 28\}$; $4 \le 7 \le 8$.

\textbf{14.} The components stay $29$, $10$ and $1$, but the \omterm{def:g10:graphs-first-look:eccentricity}{radius} becomes $5$ and the \omterm{def:g10:graphs-first-look:eccentricity}{diameter} $9$.

\textbf{15.} Without $1$--$28$, the two halves of the component are joined only through $5$--$37$; the people at
the middle of that new route, $5$ and $37$, become the closest to everyone.

\textbf{16.} Without $5$--$37$, the shortest \omterm{def:g10:graphs-first-look:path}{paths} still go through $1$--$28$, which is shorter: no \omterm{def:g10:graphs-first-look:distance}{distance} that
matters changes. Without $1$--$28$, every \omterm{def:g10:graphs-first-look:path}{path} between the two halves makes the detour through $5$--$37$.

\textbf{17.} Persons $18$ and $20$ become a component of their own: the largest component has $27$ people, with the
same \omterm{def:g10:graphs-first-look:eccentricity}{radius} and \omterm{def:g10:graphs-first-look:eccentricity}{diameter}.

\textbf{18.} Person $1$ or person $28$: the news reaches all $29$ people in $4$ days.

\textbf{19.} The largest friendship component of Brenmoor has $29$ people, \omterm{def:g10:graphs-first-look:eccentricity}{diameter} $7$, \omterm{def:g10:graphs-first-look:eccentricity}{radius} $4$ and \omterm{def:g10:graphs-first-look:eccentricity}{centre}
$\{1, 28\}$. If the friendship between the two central people, $1$ and $28$, ends, the component stays in one piece
but its \omterm{def:g10:graphs-first-look:eccentricity}{diameter} grows from $7$ to $9$.
\end{solution}
